\section{Related Work}
\label{sec:related}

\paragraph{Knowledge distillation}
Logit distillation trains a student on a teacher's softened outputs~\citep{hinton2015distilling}.
Relational KD (RKD) instead aligns pairwise distances and angles of the embeddings~\citep{park2019relational}, and contrastive distillation pursues the same idea with a different objective~\citep{tian2020contrastive}.
A trained network's predictions are determined by inner products of its last-layer features, and this feature kernel is a natural object to distill~\citep{he2022feature}.
Our kernel-alignment loss belongs to this family. It matches batch Gram matrices of teacher and student after projecting both to eight angles, and its classical twin replaces only the kernel.

\paragraph{Quantum kernels and their limits}
The projected quantum kernel~\citep{huang2021power} compares reduced states of an embedded system instead of the full state.
Whether quantum kernels help on classical data is contested.
Kernel values can concentrate exponentially, and entanglement, noise and expressivity are among the identified sources~\citep{thanasilp2024exponential}.
Tuned fidelity kernels show no advantage on classical data~\citep{slattery2023numerical}. An advantage can also vanish under noise and finite data~\citep{wang2021towards}, and bandwidth-tuned quantum kernels have been argued to behave classically~\citep{florezablan2025similarity}.
What a quantum kernel can learn depends on the inductive bias of its embedding~\citep{kubler2021inductive}, and a large benchmark studies which design choices make such kernels effective~\citep{schnabel2025scrutiny}.
These results shape our design. We pair the quantum kernel with a classical RBF kernel on the same inputs, and one quantum variant receives the same median-heuristic scale as its twin. The comparison can then separate the kernel from its calibration.

\paragraph{Quantum-assisted distillation}
QRKD~\citep{liu2025qrkd} extends RKD with quantum-kernel values. It evaluates CNNs on MNIST and CIFAR-10 and GPT-2 on three language benchmarks, and reports consistent gains over classical RKD while all deployed models stay classical.
Classical teachers have also been distilled into quantum students~\citep{li2024hybrid,hasan2023bridging}. A variational circuit has served as a soft-label teacher for a classical student~\citep{naqvi2026structured}, and quantum-inspired compression has been combined with distillation for resource-constrained deployment~\citep{kumar2026quantuminspired}.
Our setting is closest to QRKD. It differs in the data (small, imbalanced waste images), the students, the loss, the controls (a classical twin and two ablations), the reporting (paired seeds) and the end point (INT8 models on a Raspberry Pi).
We do not claim that those results fail in their own settings.

\paragraph{Edge deployment}
Integer-arithmetic-only inference~\citep{jacob2018quantization} underlies the 8-bit deployment used here, and we apply post-training quantization to finished students.
Edge-oriented studies examine accuracy--efficiency trade-offs of distillation for CNNs and vision transformers~\citep{violos2024towards}, evaluate distillation and compression on edge-class devices~\citep{paparounas2024evaluating}, and split distilled networks between device and server~\citep{matsubara2020head}.
We measure a single compact INT8 model on the device itself and report how its predictions depend on the kernels that execute it.
